\documentclass{amsart}
% For dealing with references we use the comment environment
\usepackage{verbatim}
\newenvironment{reference}{\comment}{\endcomment}
%\newenvironment{reference}{}{}
\newenvironment{slogan}{\comment}{\endcomment}
\newenvironment{history}{\comment}{\endcomment}

% For commutative diagrams we use Xy-pic
\usepackage[all]{xy}

% We use 2cell for 2-commutative diagrams.
\xyoption{2cell}
\UseAllTwocells

% We use multicol for the list of chapters between chapters
\usepackage{multicol}

% This is generally recommended for better output
\usepackage{lmodern}
\usepackage[T1]{fontenc}

% For cross-file-references

% Package for hypertext links:
\usepackage{hyperref}

% For any local file, say "hello.tex" you want to link to please

% Theorem environments.
%
\theoremstyle{plain}
\newtheorem{theorem}[subsection]{Theorem}
\newtheorem{proposition}[subsection]{Proposition}
\newtheorem{lemma}[subsection]{Lemma}

\theoremstyle{definition}
\newtheorem{definition}[subsection]{Definition}
\newtheorem{example}[subsection]{Example}
\newtheorem{exercise}[subsection]{Exercise}
\newtheorem{situation}[subsection]{Situation}

\theoremstyle{remark}
\newtheorem{remark}[subsection]{Remark}
\newtheorem{remarks}[subsection]{Remarks}

\numberwithin{equation}{subsection}

% Macros
%
\def\lim{\mathop{\mathrm{lim}}\nolimits}
\def\colim{\mathop{\mathrm{colim}}\nolimits}
\def\Spec{\mathop{\mathrm{Spec}}}
\def\Hom{\mathop{\mathrm{Hom}}\nolimits}
\def\Ext{\mathop{\mathrm{Ext}}\nolimits}
\def\SheafHom{\mathop{\mathcal{H}\!\mathit{om}}\nolimits}
\def\SheafExt{\mathop{\mathcal{E}\!\mathit{xt}}\nolimits}
\def\Sch{\mathit{Sch}}
\def\Mor{\mathop{\mathrm{Mor}}\nolimits}
\def\Ob{\mathop{\mathrm{Ob}}\nolimits}
\def\Sh{\mathop{\mathit{Sh}}\nolimits}
\def\NL{\mathop{N\!L}\nolimits}
\def\CH{\mathop{\mathrm{CH}}\nolimits}
\def\proetale{{pro\text{-}\acute{e}tale}}
\def\etale{{\acute{e}tale}}
\def\QCoh{\mathit{QCoh}}
\def\Ker{\mathop{\mathrm{Ker}}}
\def\Im{\mathop{\mathrm{Im}}}
\def\Coker{\mathop{\mathrm{Coker}}}
\def\Coim{\mathop{\mathrm{Coim}}}

% Boxtimes
%
\DeclareMathSymbol{\boxtimes}{\mathbin}{AMSa}{"02}

%
% Macros for moduli stacks/spaces
%
\def\QCohstack{\mathcal{QC}\!\mathit{oh}}
\def\Cohstack{\mathcal{C}\!\mathit{oh}}
\def\Spacesstack{\mathcal{S}\!\mathit{paces}}
\def\Quotfunctor{\mathrm{Quot}}
\def\Hilbfunctor{\mathrm{Hilb}}
\def\Curvesstack{\mathcal{C}\!\mathit{urves}}
\def\Polarizedstack{\mathcal{P}\!\mathit{olarized}}
\def\Complexesstack{\mathcal{C}\!\mathit{omplexes}}
% \Pic is the operator that assigns to X its picard group, usage \Pic(X)
% \Picardstack_{X/B} denotes the Picard stack of X over B
% \Picardfunctor_{X/B} denotes the Picard functor of X over B
\def\Pic{\mathop{\mathrm{Pic}}\nolimits}
\def\Picardstack{\mathcal{P}\!\mathit{ic}}
\def\Picardfunctor{\mathrm{Pic}}
\def\Deformationcategory{\mathcal{D}\!\mathit{ef}}

\setlength{\emergencystretch}{3em}
\begin{document}
\title[Stacks Project, Tag 0FD5]{588 --- High powers of an ample line bundle embed the complement of a closed subscheme through sections vanishing on that subscheme}
\maketitle
\noindent Copyright (C) 2005--2025 Johan de Jong. Stacks Project authors. Source: \href{https://stacks.math.columbia.edu/tag/0FD5}{Tag 0FD5}.

\noindent Editorial difficulty note (not author text). Difficulty H1: Generation of the ideal twist alone does not establish an immersion of the complement. The proof multiplies a vanishing section by a very-ample system, compares the resulting affine-chart coordinate rings and supplies a fibre-product algebra argument showing that the enlarged section system still defines an immersion..

\noindent Editorial provenance note (not author text). History: excerpt from revision \texttt{a0444\allowbreak 6e57e\allowbreak c1fbc\allowbreak 252a8\allowbreak 71afc\allowbreak ec775\allowbreak 2fb28\allowbreak 07b14}, varieties.tex, lines 10820--10906. Reference presentation was adapted; original-author context and core support blocks were retained or added. The mathematical wording is unchanged. Licensed under GNU FDL 1.2 or later; no invariant sections or cover texts. See ../licenses/COPYING. Context blocks reproduce author definitions and supporting results; they are not additional counted problems.

\begin{lemma}
\label{lemma-generate-over-complement}
Let $k$ be a field. Let $X$ be a proper scheme over $k$. Let $\mathcal{L}$
be an ample invertible $\mathcal{O}_X$-module. Let $Z \subset X$ be a
closed subscheme. Then there exists an integer $n_0$ such that for all
$n \geq n_0$ the kernel $V_n$ of
$\Gamma(X, \mathcal{L}^{\otimes n}) \to \Gamma(Z, \mathcal{L}^{\otimes n}|_Z)$
generates $\mathcal{L}^{\otimes n}|_{X \setminus Z}$ and
the canonical morphism
$$
X \setminus Z \longrightarrow \mathbf{P}(V_n)
$$
is an immersion of schemes over $k$.
\end{lemma}

\begin{proof}
Let $\mathcal{I} \subset \mathcal{O}_X$ be the quasi-coherent ideal sheaf of
$Z$. Observe that via the inclusion
$\mathcal{I} \otimes_{\mathcal{O}_X} \mathcal{L}^{\otimes n} \subset
\mathcal{L}^{\otimes n}$ we have
$V_n = \Gamma(X, \mathcal{I} \otimes_{\mathcal{O}_X} \mathcal{L}^{\otimes n})$.
Choose $n_1$ such that for $n \geq n_1$ the sheaf
$\mathcal{I} \otimes \mathcal{L}^{\otimes n}$ is globally generated, see
Properties, Proposition \href{https://stacks.math.columbia.edu/tag/01Q3}{proposition characterize ample (Tag 01Q3)}.
It follows that $V_n$ generates $\mathcal{L}^{\otimes n}|_{X \setminus Z}$
for $n \geq n_1$.

\medskip\noindent
For $n \geq n_1$ denote
$\psi_n : V_n \to
\Gamma(X \setminus Z, \mathcal{L}^{\otimes n}|_{X \setminus Z})$
the restriction map. We get a canonical morphism
$$
\varphi = \varphi_{\mathcal{L}^{\otimes n}|_{X \setminus Z}, \psi_n} :
X \setminus Z
\longrightarrow
\mathbf{P}(V_n)
$$
by Constructions, Example \href{https://stacks.math.columbia.edu/tag/0FCY}{example projective space (Tag 0FCY)}.
Choose $n_2$ such that for all $n \geq n_2$ the invertible sheaf
$\mathcal{L}^{\otimes n}$ is very ample on $X$.
We claim that $n_0 = n_1 + n_2$ works.

\medskip\noindent
Proof of the claim. Say $n \geq n_0$ and write $n = n_1 + n'$.
For $x \in X \setminus Z$ we can choose $s_1 \in V_1$ not
vanishing at $x$. Set $V' = \Gamma(X, \mathcal{L}^{\otimes n'})$.
By our choice of $n$ and $n'$ we see that the corresponding morphism
$\varphi' : X \to \mathbf{P}(V')$ is a closed immersion. Thus if we choose
$s' \in \Gamma(X, \mathcal{L}^{\otimes n'})$ not vanishing at $x$,
then $X_{s'} = (\varphi')^{-1}(D_+(s'))$ (see
Constructions, Lemma \href{https://stacks.math.columbia.edu/tag/01NK}{lemma invertible map into proj (Tag 01NK)})
is affine and $X_{s'} \to D_+(s')$ is a closed immersion.
Then $s = s_1 \otimes s' \in V_n$ does not vanish at $x$.
If $D_+(s) \subset \mathbf{P}(V_n)$ denotes the
corresponding open affine space of our projective space, then
$\varphi^{-1}(D_+(s)) = X_s \subset X \setminus Z$ (see reference above).
The open $X_s = X_{s'} \cap X_{s_1}$ is affine, see
Properties, Lemma \href{https://stacks.math.columbia.edu/tag/01PV}{lemma affine cap s open (Tag 01PV)}.
Consider the ring map
$$
\text{Sym}(V)_{(s)} \longrightarrow \mathcal{O}_X(X_s)
$$
defining the morphism $X_s \to D_+(s)$. Because $X_{s'} \to D_+(s')$
is a closed immersion, the images of the elements
$$
\frac{s_1 \otimes t'}{s_1 \otimes s'}
$$
where $t' \in V'$ generate the image of
$\mathcal{O}_X(X_{s'}) \to \mathcal{O}_X(X_s)$.
Since $X_s \to X_{s'}$ is an open immersion,
this implies that $X_s \to D_+(s)$ is an immersion of affine schemes
(see below). Thus $\varphi_n$ is an immersion by
Morphisms, Lemma \href{https://stacks.math.columbia.edu/tag/0FCZ}{lemma check immersion (Tag 0FCZ)}.

\medskip\noindent
Let $a : A' \to A$ and $c : B \to A$ be ring maps such that
$\Spec(a)$ is an immersion and $\Im(a) \subset \Im(c)$.
Set $B' = A' \times_A B$ with projections $b : B' \to B$
and $c' : B' \to A'$. By assumption $c'$
is surjective and hence $\Spec(c')$ is a closed immersion.
Whence $\Spec(c') \circ \Spec(a)$ is an immersion
(Schemes, Lemma \href{https://stacks.math.columbia.edu/tag/02V0}{lemma composition immersion (Tag 02V0)}).
Then $\Spec(c)$ has to be an immersion because it factors
the immersion $\Spec(c') \circ \Spec(a) = \Spec(b) \circ \Spec(c)$, see
Morphisms, Lemma \href{https://stacks.math.columbia.edu/tag/07RK}{lemma immersion permanence (Tag 07RK)}.
\end{proof}
\end{document}
